\chapter{数式処理システム入門}\label{quick2}

数式処理システムとは，シンボリックな計算をサポートするプログラミング言語のことをいう．
数式処理システムは，未束縛の変数をシンボルとして扱う．
数式処理システムを使うと，たとえば，$x + x \rightarrow 2x$や$(x + y)^2 \rightarrow x^2 + 2 x y + y^2$のような計算ができる．
Egisonもこのようなシンボリックな計算をサポートしている．


\section{未定義変数 = シンボル}\label{quick-symbol}

数式処理システムは，未束縛の変数をシンボルとして扱う．
そのため，未定義の変数をプログラム中で使ってもエラーにならない．
\begin{lstlisting}[language=egison,numbers=none]
x -- x
\end{lstlisting}

数式処理システムとしての機能をもつEgisonには，シンボル同士の足し算や掛け算が組み込みで定義されている．
そのため，たとえば，$x + x \rightarrow 2x$や$(x + y)^2 \rightarrow x^2 + 2 x y + y^2$のような計算ができる．
\begin{lstlisting}[language=egison,numbers=none]
x + x -- 2x
(x + y)^2 -- x^2 + 2 * y * x + y^2
\end{lstlisting}

Egisonは数式を自動で積和標準形に展開する．
積和標準形とは，掛け算が内側に，足し算が外側になるような形の数式のことである．
そのため，$(x + y)^2$は$x^2 + 2 x y + y^2$のように展開される．

\section{数式の簡約}\label{quick2-simplify}

$sin^2\theta + cos^2\theta = 1$のような数式を紙の上で扱うときにおこなう簡約が，Egisonには実装されている．
\begin{lstlisting}[language=egison,numbers=none]
(sin `$\theta$`)^2 + (cos `$\theta$`)^2 -- 1
\end{lstlisting}
このような簡約のための機能は，Sweet Egison(Egisonのパターンマッチ機能を提供するHaskellライブラリ)を使ってHaskellにより言語に組み込みで実装されている．

三角関数以外にも，複素数や平方根の簡約ルールも実装されている．
\begin{lstlisting}[language=egison,numbers=none]
i^2 -- -1
(sqrt x)^2 -- x
\end{lstlisting}

\todo{簡約ルール一覧}

